%% kapa: exam=1 round=25 year=2014 period=1 form=A subject=civil_law subject_ko=민법 seq=6 number=6
%% source: Q-Net 감정평가사 기출문제 (cst00309) artlSeq=5101207 file=extracted/1008/5101207/2014년 제25회 감정평가사 1차시험 기출문제/1교시(A형).hwp
%% conversion: kordoc 4.24.51 -> tools/gen1.py (auto); review=proofread-sonnet; flags=none
\begin{question}{6}
실종선고에 관한 설명으로 옳지 않은 것은? (다툼이 있으면 판례에 의함)
\begin{choices}
  \item 선박침몰로 인한 실종기간은 1년이고, 그 기간은 선박이 침몰한 때부터 기산한다.
  \item 실종선고가 취소되지 않는 한, 실종선고의 효과가 반증을 통하여 번복되는 것은 아니다.
  \item 법원이 실종선고 및 그 취소를 할 때에는 반드시 공시최고의 절차를 거쳐야 한다.
  \item 실종선고를 받은 자는 실종기간이 만료한 때에 사망한 것으로 본다.
  \item 실종선고에 의한 사망의 효과는 실종자의 종래의 주소나 거소를 중심으로 하는 사법적 법률관계에 국한된다.
\end{choices}
\end{question}
\answer{3}
